\slcol{अर्जुन उवाच —\\
\Index{किं तद्ब्रह्म किमध्यात्मं} किं कर्म पुरुषोत्तम । \\
अधिभूतं च किं प्रोक्तमधिदैवं किमुच्यते ॥ ८.१ ॥ }
\translate{arjuna uvāca —\\
kiṁ tadbrahma kimadhyātmaṁ kiṁ karma puruṣōttama | \\
adhibhūtaṁ ca kiṁ prōktamadhidaivaṁ kimucyatē || 8.1 || }
\cquote{Arjuna said: 
O \emph{Puruṣottama} (Supreme being, Krishna), what is \emph{Brahman}, what is the Self (\emph{adhyātma}), and what is action? What is described as the material principle (\emph{adhibhūtam}), and what is known as the divine principle (\emph{adhidaivam})?}

\slcol{\Index{अधियज्ञः कथं} कोऽत्र देहेऽस्मिन्मधुसूदन ।\\ 
प्रयाणकाले च कथं ज्ञेयोऽसि नियतात्मभिः ॥ ८.२ ॥ }
\translate{adhiyajñaḥ kathaṁ kō:'tra dēhē:'sminmadhusūdana | \\
prayāṇakālē ca kathaṁ jñēyō:'si niyatātmabhiḥ || 8.2 || }
\cquote{O Madhusūdana (Destroyer of Madhu, Krishna), who is the presiding principle of sacrifice (\emph{adhiyajña}) here in this body, and how are You to be known at the time of death by the self-controlled?}

\newpage
\begin{mananam}{\mananamfont {Mananam Śloka - 1, 2}}
\normalfont\mananamtext Have I ever delved into inquiries about the nature of the universe, its origin, and the forces that sustain it? Or have my curiosities been quenched by the answers handed down to me, leading me to now consider such inquiries impractical and a waste of time?\\
What about inquiries into the nature of life? Do I still aspire to find deeper meaning and purpose in life? What are my notions and beliefs about the afterlife and reincarnation?
\end{mananam}
\sreflect
\begin{inspiration}{\mananamfont Inspiration}
\normalfont\mananamtext The starting point for all discoveries comes from asking the right questions, whether in the field of science or spirituality. Engaging in inquiries and probing deeper is also the art of contemplation that philosophers and sages have practiced. The Buddha’s quest for truth began with an inquiry into the general nature of suffering. Ramana Maharshi’s realization came about through the inquiry ‘Who am I?’\\
The value of studying the scriptures is not to find a set of ready answers, but to learn how to ask the right questions that lead to personal realization. We must keep the lamp of inquiry aflame with the urgency of finding answers in this very life!
\end{inspiration}
\newpage

\slcol{श्रीभगवानुवाच —\\
\Index{अक्षरं ब्रह्म परमं} स्वभावोऽध्यात्ममुच्यते । \\
भूतभावोद्भवकरो विसर्गः कर्मसंज्ञितः ॥ ८.३ ॥ }
\translate{śrībhagavānuvāca —\\
akṣaraṁ brahma paramaṁ svabhāvō:'dhyātmamucyatē | \\
bhūtabhāvōdbhavakarō visargaḥ karmasaṁjñitaḥ || 8.3 || }
\cquote{Sri Bhagavan said:
The imperishable \emph{Brahman} is supreme. The essential nature of a being is called the Self. The creative force that brings beings into existence is known as \emph{karma}.}

\slcol{\Index{अधिभूतं क्षरो भावः} पुरुषश्चाधिदैवतम् । \\
अधियज्ञोऽहमेवात्र देहे देहभृतां वर ॥ ८.४ ॥ }
\translate{adhibhūtaṁ kṣarō bhāvaḥ puruṣaścādhidaivatam | \\
adhiyajñō:'hamēvātra dēhē dēhabhr̥tāṁ vara || 8.4 || }
\cquote{The material principle (\emph{adhibhūtam}) is My perishable nature; the divine principle (\emph{adhidaivam}) is the \emph{Puruṣa} (Conscious being or Consciousness from \emph{Sāṅkhya} philosophy). And I alone am the \emph{adhiyajña} here in the body, O best of embodied beings (Arjuna).}

\slcol{\Index{अन्तकाले च मामेव} स्मरन्मुक्त्वा कलेबरम् । \\
यः प्रयाति स मद्भावं याति नास्त्यत्र संशयः ॥ ८.५ ॥ }
\translate{antakālē ca māmēva smaranmuktvā kalēbaram | \\
yaḥ prayāti sa madbhāvaṁ yāti nāstyatra saṁśayaḥ || 8.5 || }
\cquote{Whoever, at the final moment, remembers Me alone and leaves the body, attains My state of being, without doubt.}

\slcol{\Index{यं यं वापि स्मरन्भावं} त्यजत्यन्ते कलेबरम् । \\
तं तमेवैति कौन्तेय सदा तद्भावभावितः ॥ ८.६ ॥ }
\translate{yaṁ yaṁ vāpi smaranbhāvaṁ tyajatyantē kalēbaram | \\
taṁ tamēvaiti kauntēya sadā tadbhāvabhāvitaḥ || 8.6 || }
\cquote{At the final moment, whatever thought or state one remembers, that he attains, O Arjuna, being shaped by constant absorption in it.}

\newpage
\begin{mananam}{\mananamfont{Mananam Śloka - 5, 6}}
\normalfont\mananamtext Have I ever considered the importance of the state of mind at the moment of death? Or am I solely focused on enjoying life and having pleasurable experiences daily? Have I experienced the connection between the nature of my thoughts prior to sleep and how they influence my dreams?\\
How can I ensure that my thoughts at the moment of leaving my body are centred on God or the Divinity within me? During difficult moments in my life, do my thoughts turn to God?
\end{mananam}
\sreflect
\begin{inspiration}{\mananamfont Inspiration}
\normalfont\mananamtext All our thoughts and efforts leave imprints in the unseen areas of the mind. These imprints deeply influence all our activities. The scriptures assert that our next life (rebirth in any realm) is determined by the most dominant and reaffirmed impressions. In fact, other scriptures point out that not just rebirth but liberation can be attained at the crucial moment of transition from this body. If our thoughts at the time of death can hold onto that One Divine essence, it is the most effective approach to liberation. However, we can only contemplate any one given desired thing at the point of death if we have practiced it throughout our life. Thus, for a \emph{sādhaka}, all life is an opportunity to practice being in the Divine presence.
\end{inspiration}
\newpage

\slcol{\Index{तस्मात्सर्वेषु कालेषु} मामनुस्मर युध्य च । \\
मय्यर्पितमनोबुद्धिर्मामेवैष्यस्यसंशयः ॥ ८.७ ॥ }
\translate{tasmātsarvēṣu kālēṣu māmanusmara yudhya ca | \\
mayyarpitamanōbuddhirmāmēvaiṣyasyasaṁśayaḥ || 8.7 || }
\cquote{Therefore, at all times remember Me and fight (perform your duty). With mind and intellect offered to Me, you shall surely come to Me, without doubt.}

\slcol{\Index{अभ्यासयोगयुक्तेन} चेतसा नान्यगामिना । \\
परमं पुरुषं दिव्यं याति पार्थानुचिन्तयन् ॥ ८.८ ॥ }
\translate{abhyāsayōgayuktēna cētasā nānyagāminā | \\
paramaṁ puruṣaṁ divyaṁ yāti pārthānucintayan || 8.8 || }
\cquote{With the mind disciplined by the regular practice of \emph{yoga}, not distracted elsewhere, one attains the supreme, divine \emph{Puruṣa} by constantly meditating upon Him, O Pārtha (Arjuna).}

\slcol{\Index{कविं पुराणमनुशासितार}मणोरणीयांसमनुस्मरेद्यः । \\
सर्वस्य धातारमचिन्त्यरूपमादित्यवर्णं तमसः परस्तात् ॥ ८.९ ॥ }
\translate{kaviṁ purāṇamanuśāsitāramaṇōraṇīyāṁsamanusmarēdyaḥ | \\
sarvasya dhātāramacintyarūpamādityavarṇaṁ \\\hspace*{0.5cm}tamasaḥ parastāt || 8.9 || }
\cquote{He who meditates upon the omniscient, primeval ruler, subtler than the subtle, the sustainer of all, of inconceivable form, radiant like the sun, beyond all darkness…}

\slcol{\Index{प्रयाणकाले मनसाचलेन} 
भक्त्या युक्तो योगबलेन चैव । \\
भ्रुवोर्मध्ये प्राणमावेश्य सम्य 
क्स तं परं पुरुषमुपैति दिव्यम् ॥८.१०॥ }
\translate{prayāṇakālē manasācalēna 
bhaktyā yuktō yōgabalēna caiva | \\
bhruvōrmadhyē prāṇamāvēśya samya \\
\hspace*{0.5cm}ksa taṁ paraṁ puruṣamupaiti divyam || 8.10 || }
\cquote{At the time of departure, with mind unwavering, devotion steady, and by the power of \emph{yoga}, fixing the life-breath between the eyebrows, one attains the supreme, divine \emph{Puruṣa.}}

\newpage
\slcol{\Index{यदक्षरं वेदविदो} वदन्ति 
विशन्ति यद्यतयो वीतरागाः । \\
यदिच्छन्तो ब्रह्मचर्यं चरन्ति 
तत्ते पदं सङ्ग्रहेण प्रवक्ष्ये ॥ ८.११ ॥ }
\translate{yadakṣaraṁ vēdavidō vadanti 
viśanti yadyatayō vītarāgāḥ | \\
yadicchantō brahmacaryaṁ caranti \\
\hspace*{0.5cm}tattē padaṁ saṅgrahēṇa pravakṣyē || 8.11 || }
\cquote{That imperishable state, which the knowers of the \emph{Veda} declare, into which the self-controlled and passionless ascetics enter, desiring which they practice celibacy—that state I shall briefly declare to you.}

\slcol{\Index{सर्वद्वाराणि संयम्य} मनो हृदि निरुध्य च । \\
मूर्ध्न्याधायात्मनः प्राणमास्थितो योगधारणाम् ॥ ८.१२ ॥ }
\translate{sarvadvārāṇi saṁyamya manō hr̥di nirudhya ca | \\
mūrdhnyādhāyātmanaḥ prāṇamāsthitō yōgadhāraṇām || 8.12 || }
\cquote{Restraining all the gates of the body (senses), fixing the mind in the heart, and directing the life-breath to the crown of the head, one becomes firmly established in the practice of \emph{yoga}.}

\slcol{\Index{ओमित्येकाक्षरं ब्रह्म} 
व्याहरन्मामनुस्मरन् । \\
यः प्रयाति त्यजन्देहं 
स याति परमां गतिम् ॥ ८.१३ ॥ }
\translate{ōmityēkākṣaraṁ brahma 
vyāharanmāmanusmaran | \\
yaḥ prayāti tyajandēhaṁ 
sa yāti paramāṁ gatim || 8.13 || }
\cquote{Uttering the single syllable \emph{Om}, which is (the symbol of) \emph{Brahman}, and remembering Me, whoever departs leaving the body attains the supreme goal.}

\slcol{\Index{अनन्यचेताः सततं} यो मां स्मरति नित्यशः । \\
तस्याहं सुलभः पार्थ नित्ययुक्तस्य योगिनः ॥ ८.१४ ॥ }
\translate{ananyacētāḥ satataṁ yō māṁ smarati nityaśaḥ | \\
tasyāhaṁ sulabhaḥ pārtha nityayuktasya yōginaḥ || 8.14 || }
\cquote{O Pārtha, I am easily attainable by the \emph{yogi} who is ever united, who remembers Me constantly with single-minded devotion.}

\newpage
\begin{mananam}{\mananamfont{Mananam Śloka - 10, 13}}
\normalfont\mananamtext Just as there are various medical sciences for physical and mental well-being, could there be a lesser-known science that reveals the path for the soul departing the physical body? Do my doubts and lack of trust about the afterlife and reincarnation deter me from undertaking any investigation on my own?\\
Have I practiced or undertaken any serious investigation to gain personal experience or understanding of whether there is a soul beyond this body-mind? Have I explored whether there is a supreme entity and its nature may be, which we conceptually refer to as “God”?
\end{mananam}
\sreflect
\begin{inspiration}{\mananamfont{Inspiration}}
\normalfont\mananamtext The Lord provides a specific technique to attain Him, as mere words may not suffice for most. This technique involves using the sacred syllable \emph{OM}, stilling the mind, and focusing attention on the spiritual energy centre in the body, located between the eyebrows. To effectively use this technique at the point of death, one must practice it during life and verify its effectiveness. Just as a warrior who practices regularly is prepared for battle, so too must we prepare ourselves through consistent practice in life so that we are prepared to face death.
\end{inspiration}
\newpage

\slcol{\Index{मामुपेत्य पुनर्जन्म} दुःखालयमशाश्वतम् । \\
नाप्नुवन्ति महात्मानः संसिद्धिं परमां गताः ॥ ८.१५ ॥ }
\translate{māmupētya punarjanma duḥkhālayamaśāśvatam | \\
nāpnuvanti mahātmānaḥ saṁsiddhiṁ paramāṁ gatāḥ || 8.15 || }
\cquote{Having attained Me, the great souls do not return to rebirth, which is the abode of sorrow, and impermanent. They have reached the supreme perfection.}

\slcol{\Index{आब्रह्मभुवनाल्लोकाः} पुनरावर्तिनोऽर्जुन । \\
मामुपेत्य तु कौन्तेय पुनर्जन्म न विद्यते ॥ ८.१६ ॥ }
\translate{ā brahmabhuvanāllōkāḥ punarāvartinō:'rjuna | \\
māmupētya tu kauntēya punarjanma na vidyatē || 8.16 || }
\cquote{O Arjuna, all worlds—even up to Brahmā’s abode—are subject to return. But those who attain Me never face rebirth again, O Kaunteya (son of Kunti, Arjuna).}

\slcol{\Index{सहस्रयुगपर्यन्तमहर्यद्ब्रह्मणो} विदुः । \\
रात्रिं युगसहस्रान्तां तेऽहोरात्रविदो जनाः ॥ ८.१७ ॥ }
\translate{sahasrayugaparyantamaharyadbrahmaṇō viduḥ | \\
rātriṁ yugasahasrāntāṁ tē:'hōrātravidō janāḥ || 8.17 || }
\cquote{Those who know day and night (cosmic time) understand that the day of Brahmā lasts for a thousand \emph{yugas} (aeons), and his night also lasts for a thousand \emph{yugas}.}

\slcol{\Index{अव्यक्ताद्व्यक्तयः सर्वाः} प्रभवन्त्यहरागमे । \\
रात्र्यागमे प्रलीयन्ते तत्रैवाव्यक्तसंज्ञके ॥ ८.१८ ॥ }
\translate{avyaktādvyaktayaḥ sarvāḥ prabhavantyaharāgamē | \\
rātryāgamē pralīyantē tatraivāvyaktasaṁjñakē || 8.18 || }
\cquote{From the unmanifest, all beings come forth at the arrival of day (of Brahmā); at the arrival of night, they dissolve into that same unmanifest.}

\slcol{\Index{भूतग्रामः स एवायं} भूत्वा भूत्वा प्रलीयते । \\
रात्र्यागमेऽवशः पार्थ प्रभवत्यहरागमे ॥ ८.१९ ॥ }
\translate{bhūtagrāmaḥ sa ēvāyaṁ bhūtvā bhūtvā pralīyatē | \\
rātryāgamē:'vaśaḥ pārtha prabhavatyaharāgamē || 8.19 || }
\cquote{O Pārtha (Arjuna), this multitude of beings repeatedly comes into existence and dissolves at Brahmā’s night; helplessly, it arises again at the dawn of his day.}

\slcol{\Index{परस्तस्मात्तु भावोऽन्योऽव्यक्तो}ऽव्यक्तात्सनातनः । \\
यः स सर्वेषु भूतेषु नश्यत्सु न विनश्यति ॥ ८.२० ॥ }
\translate{parastasmāttu bhāvō:'nyō:'vyaktō:'vyaktātsanātanaḥ | \\
yaḥ sa sarvēṣu bhūtēṣu naśyatsu na vinaśyati || 8.20 || }
\cquote{Yet beyond this unmanifest lies another, eternal unmanifest. Though all beings perish, that reality does not perish.}

\slcol{\Index{अव्यक्तोऽक्षर इत्युक्तस्तमाहुः} परमां गतिम् । \\
यं प्राप्य न निवर्तन्ते तद्धाम परमं मम ॥ ८.२१ ॥ }
\translate{avyaktō:'kṣara ityuktastamāhuḥ paramāṁ gatim | \\
yaṁ prāpya na nivartantē taddhāma paramaṁ mama || 8.21 || }
\cquote{That unmanifest, imperishable reality is spoken of as the supreme goal. Having attained it, one never returns; it is My supreme abode.}


\slcol{\Index{पुरुषः स परः} पार्थ भक्त्या लभ्यस्त्वनन्यया । \\
यस्यान्तःस्थानि भूतानि येन सर्वमिदं ततम् ॥ ८.२२ ॥ }
\translate{puruṣaḥ sa paraḥ pārtha bhaktyā labhyastvananyayā | \\
yasyāntaḥsthāni bhūtāni yēna sarvamidaṁ tatam || 8.22 || }
\cquote{O Pārtha (Arjuna), the Supreme Person is attained through unwavering devotion alone. All beings abide in Him, and He pervades the entire universe.}

\slcol{\Index{यत्र काले त्वनावृत्तिमावृत्तिं} चैव योगिनः । \\
प्रयाता यान्ति तं कालं वक्ष्यामि भरतर्षभ ॥ ८.२३ ॥ }
\translate{yatra kālē tvanāvr̥ttimāvr̥ttiṁ caiva yōginaḥ | \\
prayātā yānti taṁ kālaṁ vakṣyāmi bharatarṣabha || 8.23 || }
\cquote{O Bharatarṣabha (best of the Bharatas), I will now explain the times of departure by which \emph{yogis} attain liberation and by which they return again (are reborn).}

\newpage
\begin{mananam}{\mananamfont{Mananam Śloka - 20, 21}}
\normalfont\mananamtext What is my understanding of the unmanifest (\emph{Avyakta}), where the seed of all creation lies? Can there be a time when the world comes to an end? Every night, when I enter deep sleep, don’t I experience the dissolution of my perception of this world? When I wake up in the morning, how does this world re-emerge for me? How am I able to recollect memories, regain self-awareness, and become aware of the world again? Was I present during deep sleep or not? Similarly, can the common experience of the world also undergo dissolution and re-emergence, as the Hindu scriptures suggest?
\end{mananam}
\sreflect
\begin{inspiration}{\mananamfont{Inspiration}}
\normalfont\mananamtext One of the greatest wonders and mysteries of the world is how an entire tree emerges from a seed. Just as the design of a tree is contained within a seed, the imprint of the entire universe lies in the Unmanifest state before creation and dissolves back into that state at dissolution. \\
The eternal Consciousness or Existence that remains beyond the unmanifest is the ultimate destination of all souls. Finding our path back to that reality leads us to the ultimate resting place. The entire science of \emph{yoga} aims to guide us back to this essential reality, freeing us from the cycles of birth and death.
\end{inspiration}
\newpage
\begin{mananam}{\mananamfont{Mananam Śloka - 22}}
\normalfont\mananamtext Why should I engage in inquiries about the unmanifest state and the Supreme beyond it? Can I not be content like the majority who follow the dogma, “Seeing is believing”?\\
How can understanding the unmanifest state and the Supreme help me know more about myself? Can the essential Source of all creation be present within this creation? What is the true power behind my abilities to see, hear, think, breathe, and digest? Is there one essential attribute common to all life, similar to how breath is essential for physical life?
\end{mananam}
\sreflect
\begin{inspiration}{\mananamfont{Inspiration}}
\normalfont\mananamtext How readily we give our love and devotion to our family and friends! How naturally we turn our loyalty towards our bosses and those in power! Yet, how little do we focus on the essential Source of all creation. Our own awareness, through which we experience and navigate this world and think our thoughts, is infused with energy and consciousness. The scriptures and sages prod us to come out of common human ignorance. They urge us to understand that essence which pervades the whole of creation, though it is not visible to our eyes and senses. 
\end{inspiration}
\newpage

\slcol{\Index{अग्निर्ज्योतिरहः शुक्लः} षण्मासा उत्तरायणम् । \\
तत्र प्रयाता गच्छन्ति ब्रह्म ब्रह्मविदो जनाः ॥ ८.२४ ॥ }
\translate{agnirjyōtirahaḥ śuklaḥ ṣaṇmāsā uttarāyaṇam | \\
tatra prayātā gacchanti brahma brahmavidō janāḥ || 8.24 || }
\cquote{Departing in the time of fire, light, day, the bright fortnight, and the six months of the sun’s northern path, the knowers of \emph{Brahman} attain \emph{Brahman}.}


\slcol{\Index{धूमो रात्रिस्तथा कृष्णः} षण्मासा दक्षिणायनम् । \\
तत्र चान्द्रमसं ज्योतिर्योगी प्राप्य निवर्तते ॥ ८.२५ ॥ }
\translate{dhūmō rātristathā kr̥ṣṇaḥ ṣaṇmāsā dakṣiṇāyanam | \\
tatra cāndramasaṁ jyōtiryōgī prāpya nivartatē || 8.25 || }
\cquote{Departing in the time of smoke, night, the dark fortnight, and the six months of the sun’s southern path, the \emph{yogi} attains the moon’s light and returns.}


\slcol{\Index{शुक्लकृष्णे गती ह्येते} जगतः शाश्वते मते । \\
एकया यात्यनावृत्तिमन्ययावर्तते पुनः ॥ ८.२६ ॥ }
\translate{śuklakr̥ṣṇē gatī hyētē jagataḥ śāśvatē matē | \\
ekayā yātyanāvr̥ttimanyayāvartatē punaḥ || 8.26 || }
\cquote{The bright and dark paths are eternal in the world’s understanding. By one (the path of light), there is no return; by the other (the path of darkness), one returns again.}

\slcol{\Index{नैते सृती पार्थ} जानन्योगी मुह्यति कश्चन । \\
तस्मात्सर्वेषु कालेषु योगयुक्तो भवार्जुन ॥ ८.२७ ॥ }
\translate{naitē sr̥tī pārtha jānanyōgī muhyati kaścana | \\
tasmātsarvēṣu kālēṣu yōgayuktō bhavārjuna || 8.27 || }
\cquote{O Pārtha, the \emph{yogis} are not deluded by these two paths. Therefore, at all times, O Arjuna, be steadfast in \emph{yoga}. }

\newpage
\slcol{\Index{वेदेषु यज्ञेषु तपःसु चैव} 
दानेषु यत्पुण्यफलं प्रदिष्टम् । \\
अत्येति तत्सर्वमिदं विदित्वा 
योगी परं स्थानमुपैति चाद्यम् ॥ ८.२८ ॥ }
\translate{vēdēṣu yajñēṣu tapaḥsu caiva 
dānēṣu yatpuṇyaphalaṁ pradiṣṭam | \\
atyēti tatsarvamidaṁ viditvā 
yōgī paraṁ sthānamupaiti cādyam || 8.28 || }
\cquote{The \emph{yogi}, having known this teaching, surpasses all the merits declared from \emph{vedic} study, sacrifices, austerities, and gifts, and attains the supreme, eternal abode.}

\newpage
\begin{mananam}{\mananamfont{Mananam Śloka - 28}}
\normalfont\mananamtext How am I using the life choices available to me in this life based on the truths of the afterlife? To what extent do I believe in and practice the \emph{vedic} injunctions related to rites, rituals, and offerings that are said to benefit my afterlife and that of my ancestors?
Do I participate in charitable activities and good deeds merely for a sense of well-being, or also with the hope that their merits might accrue to me beyond this life? 
What knowledge or understanding have I acquired, or what practices am I undertaking now, which could aid me at the point of death? Are there any scripturally based teachings or practices that might help me prepare for the moment of death and beyond? 
\end{mananam}
\sreflect
\begin{inspiration}{\mananamfont{Inspiration}}
\normalfont\mananamtext Just as in life, where we are always presented with at least two options, so at the point of death, every departing soul has two paths to choose from: one that leads to a higher plane of existence and another that brings the majority back to Earth or to lower planes of existence, based on their stock of merits and sins. The \emph{vedas} provide various rituals and sacrifices to aid in following the former path. However, the goal of a \emph{yogi} is neither of these paths, as any choice therein is still bound by limitation. Through realization and abidance in Truth, a \emph{yogi} attains the Supreme state of \emph{Akṣara Brahma}, which is beyond all creation and beyond one’s stock of virtues and sins. Having found that Completeness, the \emph{yogi} seeks nothing else. 
\end{inspiration}
\newpage

\begin{center}
\devfont ॐ तत्सत्।  इति श्रिमद्भगवद्गीतासूपनिष्त्सु ब्रह्मविद्यायां योगशास्रे\\
श्रीकृष्णार्जुनसंवादे अक्षरब्रह्मयोगो नाम अष्टमोऽध्यायः ॥\\
\chapEndSlokaTranslate{ōṁ tatsaditi śrīmadbhagavadgītāsūpaniṣatsu brahmavidyāyāṁ yōgaśāstrē śrīkṛṣṇārjunasaṁvādē akṣarabrahmayōgō nāma aṣṭamō'dhyāyaḥ||}
\end{center}